% ======================== 第四节 发行人基本情况 ==============================
% 基本信息 → 设立沿革与折股 → 发行前后股本结构 → 控股股东/实际控制人
% → 股权结构图（TikZ 自绘，无外部图片依赖）→ 子公司 → 董事、高级管理人员
% 与核心技术人员 → 员工。治理结构按新《公司法》：不设监事会，审计委员会
% 行使监事会职权；职工人数超过三百人，董事会设职工代表董事。
\gssection{发行人基本情况}

\subsection{发行人基本情况}

\begin{gstab}
\begin{longtable}{|L{35mm}|L{110mm}|}
\hline
中文名称 & \gsCompany \\ \hline
英文名称 & \gsCompanyEn \\ \hline
注册资本 & 36,000.00万元 \\ \hline
法定代表人 & 林×× \\ \hline
有限公司成立日期 & 2021年3月18日 \\ \hline
股份公司设立日期 & 2022年9月28日 \\ \hline
住所 & \gsRegAddr \\ \hline
邮政编码 & ×××××× \\ \hline
电话／传真 & ××××-××××××××\ ／\ ××××-×××××××× \\ \hline
互联网网址 & www.example.com \\ \hline
电子信箱 & ir@example.com \\ \hline
负责信息披露和投资者关系的部门 & 董事会办公室（联系人：林×雪） \\ \hline
统一社会信用代码 & 91××××××××××××××××\\ \hline
\end{longtable}
\end{gstab}

\subsection{发行人设立情况及报告期内的股本变化}

\subsubsection{发行人的设立方式}

公司前身示例科技有限公司于2021年3月18日设立，设立时注册资本1,000.00万元。
设立后，示例有限先后引入示例创投、示例产投等外部机构投资者增资，并设立员工
持股平台示例聚力实施股权激励；截至整体变更基准日2022年6月30日，示例有限经
审计的净资产为45,682.35万元。

2022年9月28日，示例有限按经审计净资产值折股整体变更为股份有限公司，以截至
2022年6月30日经审计的净资产45,682.35万元按1.2690:1的比例折合股本36,000.00
万股，余额计入资本公积。本次整体变更经示例会计师验资（示例会所验字〔2022〕
第××号），并已办理工商变更登记。

\subsubsection{报告期内的股本变化}

报告期内，公司股本总额未发生变化，未发生增资扩股或股份回购事项；股东之间的
少量股份转让不影响公司控制权结构。

\subsubsection{公司设立以来的主要发展历程}

% 发展历程时间轴（式样对照 2026 年申报稿的产品/发展演进图：横向箭头时间轴，
% 里程碑上下交错标注）
\begin{center}
\begin{tikzpicture}[x=1mm, y=1mm,
  ms/.style={font=\zihao{-5}, align=center},
  yr/.style={font=\heiti\zihao{5}, text=gsbar}]
\draw[gsbar, line width=1.2pt, -{Stealth[length=3mm]}] (0,0) -- (146,0);
% 轴上里程碑（奇数位在上）
\foreach \x/\yy/\txt in {
  16/2021.3/{示例有限设立，启动\\基础模型自主研发},
  74/2023.6/{开放平台（MaaS）上线，\\完成生成式AI服务备案},
  132/2025.12/{智能体平台发布，全年\\营业收入超13亿元}}{
  \fill[gsbar] (\x,0) circle (1.6);
  \node[yr, anchor=south] at (\x,3) {\yy};
  \node[ms, anchor=south] at (\x,9) {\txt};
}
% 轴下里程碑（偶数位在下）
\foreach \x/\yy/\txt in {
  45/2022.9/{整体变更设立股份公司，\\首代基础模型发布},
  103/2024.8/{多模态模型发布，金融、\\政务标杆项目规模化落地}}{
  \fill[gsbar] (\x,0) circle (1.6);
  \node[yr, anchor=north] at (\x,-3) {\yy};
  \node[ms, anchor=north] at (\x,-9) {\txt};
}
\end{tikzpicture}
\end{center}
\gsfigcap{公司发展历程}

\subsection{本次发行前后的股本结构}

本次发行前，公司总股本为36,000.00万股；本次拟公开发行新股12,000.00万股，发
行后总股本为48,000.00万股。本次发行前后的股本结构如下：

\begin{gstab}
\begin{longtable}{|L{38mm}|R{26mm}|R{20mm}|R{26mm}|R{20mm}|}
\hline
\thA{\multirow{2}{*}{股东名称}} & \multicolumn{2}{c|}{\bfseries 本次发行前} & \multicolumn{2}{c|}{\bfseries 本次发行后} \\ \cline{2-5}
\thA{} & \thB{持股数（万股）} & \thB{持股比例} & \thB{持股数（万股）} & \thB{持股比例} \\ \hline
\endfirsthead
\hline
\thA{\multirow{2}{*}{股东名称}} & \multicolumn{2}{c|}{\bfseries 本次发行前} & \multicolumn{2}{c|}{\bfseries 本次发行后} \\ \cline{2-5}
\thA{} & \thB{持股数（万股）} & \thB{持股比例} & \thB{持股数（万股）} & \thB{持股比例} \\ \hline
\endhead
示例控股 & 16,200.00 & 45.00\% & 16,200.00 & 33.75\% \\ \hline
林×× & 3,600.00 & 10.00\% & 3,600.00 & 7.50\% \\ \hline
示例创投 & 4,320.00 & 12.00\% & 4,320.00 & 9.00\% \\ \hline
示例产投 & 3,240.00 & 9.00\% & 3,240.00 & 6.75\% \\ \hline
示例聚力 & 2,880.00 & 8.00\% & 2,880.00 & 6.00\% \\ \hline
其他股东 & 5,760.00 & 16.00\% & 5,760.00 & 12.00\% \\ \hline
本次发行新股 & —— & —— & 12,000.00 & 25.00\% \\ \hline
\multicolumn{1}{|c|}{\bfseries 合计} & {\bfseries 36,000.00} & {\bfseries 100.00\%} & {\bfseries 48,000.00} & {\bfseries 100.00\%} \\ \hline
\end{longtable}
\end{gstab}
\tabnote{其他股东为公司整体变更前的历史股东，合计18名，均不属于控股股东、实
际控制人的一致行动人。}

\subsection{控股股东及实际控制人情况}

\subsubsection{控股股东}

公司控股股东为示例控股有限公司，成立于2021年2月，注册资本5,000.00万元，
林××持有其100\%股权。示例控股为持股型公司，除持有发行人股份外未开展其他实
质经营业务。本次发行前，示例控股直接持有公司16,200.00万股股份，占总股本的
45.00\%。

\subsubsection{实际控制人}

公司实际控制人为林××，中国国籍，无境外永久居留权。林××通过示例控股控制
公司45.00\%的表决权，直接持有公司10.00\%的股份，并通过担任员工持股平台示例
聚力的执行事务合伙人控制公司8.00\%的表决权，合计控制公司63.00\%的表决权。
最近两年内，公司控股股东及实际控制人未发生变更。

\subsection{发行人股权结构及对其他企业的重要权益投资}

\subsubsection{本次发行前的股权结构}

截至本招股说明书签署日，公司股权结构如下图所示：

\begin{center}
\begin{tikzpicture}[
  box/.style  = {draw, semithick, minimum height=8mm, inner sep=3pt,
                 font=\zihao{-5}, align=center},
  arr/.style  = {-{Stealth[length=2mm]}, semithick},
  lab/.style  = {font=\zihao{-5}, inner sep=1pt},
]
\node[box] (zhang)   at (2.70, 4.7) {林××};
\node[box, minimum width=26mm] (konggu)    at ( 1.300, 3.0) {示例控股};
\node[box, minimum width=26mm] (juli)      at ( 4.125, 3.0) {示例聚力};
\node[box, minimum width=26mm] (chuangtou) at ( 6.950, 3.0) {示例创投};
\node[box, minimum width=26mm] (chantou)   at ( 9.775, 3.0) {示例产投};
\node[box, minimum width=26mm] (qita)      at (12.600, 3.0) {其他股东};
\node[box, minimum width=88mm] (issuer) at (6.95, 1.1) {\gsCompany};
\node[box, minimum width=34mm] (jingmi)   at (4.60, -0.6) {示例智算};
\node[box, minimum width=34mm] (ruanjian) at (9.30, -0.6) {示例数据};
\draw[arr] (zhang.south) -- +(0,-0.45) -| (konggu.north);
\node[lab, left]  at (1.30, 3.62) {100\%};
\draw[arr] (zhang.south) -- +(0,-0.45) -| (juli.north);
\node[lab, right] at (4.125, 3.62) {执行事务合伙人};
\draw[arr] (konggu.south) -- +(0,-0.55) -| ([xshift=-34mm]issuer.north);
\node[lab, left]  at (1.30, 2.32) {45.00\%};
\draw[arr] (juli.south) -- (juli.south |- issuer.north)
  node[lab, midway, left] {8.00\%};
\draw[arr] (chuangtou.south) -- (chuangtou.south |- issuer.north)
  node[lab, midway, right] {12.00\%};
\draw[arr] (chantou.south) -- (chantou.south |- issuer.north)
  node[lab, midway, right] {9.00\%};
\draw[arr] (qita.south) -- +(0,-0.55) -| ([xshift=34mm]issuer.north);
\node[lab, right] at (12.60, 2.32) {16.00\%};
\draw[arr] (zhang.west) -- ++(-2.7,0) |- (issuer.west);
\node[lab, right] at (-0.52, 1.85) {10.00\%};
\draw[arr] (issuer.south -| jingmi.north) -- (jingmi.north)
  node[lab, midway, left] {100\%};
\draw[arr] (issuer.south -| ruanjian.north) -- (ruanjian.north)
  node[lab, midway, right] {100\%};
\end{tikzpicture}
\end{center}

\subsubsection{发行人控股子公司情况}

截至本招股说明书签署日，公司拥有2家全资子公司，无参股公司：

\begin{gstab}
\begin{longtable}{|L{34mm}|C{22mm}|C{20mm}|C{16mm}|L{34mm}|}
\hline
\thA{公司名称} & \thB{注册资本（万元）} & \thB{持股比例} & \thB{成立时间} & \thB{主营业务} \\ \hline
\endfirsthead
\hline
\thA{公司名称} & \thB{注册资本（万元）} & \thB{持股比例} & \thB{成立时间} & \thB{主营业务} \\ \hline
\endhead
示例智算科技有限公司 & 5,000.00 & 100\% & 2022年5月 & 智算集群建设运维与算力服务 \\ \hline
示例数据科技有限公司 & 2,000.00 & 100\% & 2023年2月 & 训练语料建设、数据处理与标注服务 \\ \hline
\end{longtable}
\end{gstab}

\subsection{董事、高级管理人员与核心技术人员}

\subsubsection{基本情况}

截至本招股说明书签署日，公司董事会由9名董事组成，其中独立董事3名、职工代表
董事1名；高级管理人员5名；核心技术人员4名。根据《公司法》及《公司章程》的
规定，公司不设监事会，由董事会审计委员会行使《公司法》规定的监事会职权。

\begin{gstab}
\begin{longtable}{|C{16mm}|L{50mm}|C{12mm}|C{18mm}|C{34mm}|}
\hline
\thA{姓名} & \thB{职务} & \thB{性别} & \thB{出生年份} & \thB{任期} \\ \hline
\endfirsthead
\hline
\thA{姓名} & \thB{职务} & \thB{性别} & \thB{出生年份} & \thB{任期} \\ \hline
\endhead
林×× & 董事长 & 男 & 2001年 & 2025.09—2028.09 \\ \hline
林×杰 & 董事、总经理 & 男 & 2002年 & 2025.09—2028.09 \\ \hline
林×桐 & 董事、副总经理 & 男 & 2001年 & 2025.09—2028.09 \\ \hline
林×阳 & 董事、财务总监 & 女 & 2003年 & 2025.09—2028.09 \\ \hline
林×雪 & 董事、董事会秘书 & 女 & 2002年 & 2025.09—2028.09 \\ \hline
王×× & 职工代表董事 & 男 & 2001年 & 2025.09—2028.09 \\ \hline
赵×× & 独立董事 & 女 & 2001年 & 2025.09—2028.09 \\ \hline
钱×× & 独立董事 & 男 & 2002年 & 2025.09—2028.09 \\ \hline
孙×× & 独立董事 & 男 & 2003年 & 2025.09—2028.09 \\ \hline
卫×× & 副总经理、核心技术人员 & 男 & 1998年 & 2025.09—2028.09 \\ \hline
\end{longtable}
\end{gstab}
\tabnote{核心技术人员为卫××、蒋××、沈××、韩××，其中蒋××、沈××、韩
××未在公司担任董事、高级管理人员职务。董事会审计委员会由独立董事赵××、
钱××、孙××组成，赵××（会计专业人士）任召集人。}

\subsubsection{主要成员简历}

\paragraph{林××，董事长}
2001年出生，中国国籍，无境外永久居留权，××大学计算机科学与技术专业。本科
在读期间即从事深度学习与大模型方向研究，2021年3月创立示例有限，历任执行董
事、总经理；现任公司董事长，兼任示例控股执行董事、示例聚力执行事务合伙人。

\paragraph{林×杰，董事、总经理}
2002年出生，中国国籍，无境外永久居留权，××大学软件工程专业。曾任××互联
网企业算法工程师；2021年6月加入公司，历任技术负责人、副总经理；现任公司董
事、总经理。

\paragraph{林×阳，董事、财务总监}
2003年出生，中国国籍，无境外永久居留权，××大学会计学专业，中国注册会计师。
曾任××会计师事务所审计人员；2023年3月加入公司，现任公司董事、财务总监。

\paragraph{林×雪，董事、董事会秘书}
2002年出生，中国国籍，无境外永久居留权，××大学法学专业。曾任××上市公司
证券事务代表；2024年1月加入公司，现任公司董事、董事会秘书。

其余董事、高级管理人员及核心技术人员简历从略，正式申报文件中应按上述格式逐
一披露。

\subsection{员工情况}

截至2025年12月31日，公司及子公司在册员工合计1,286人，全部员工均已签订劳动
合同，公司依法为员工缴纳社会保险及住房公积金。员工构成如下：

\begin{gstab}
\begin{longtable}{|L{46mm}|R{28mm}|R{28mm}|}
\hline
\thA{专业构成} & \thB{人数（人）} & \thB{占比} \\ \hline
\endfirsthead
\hline
\thA{专业构成} & \thB{人数（人）} & \thB{占比} \\ \hline
\endhead
研发与技术人员 & 726 & 56.45\% \\ \hline
解决方案交付与运营人员 & 268 & 20.84\% \\ \hline
销售与市场人员 & 158 & 12.29\% \\ \hline
管理及职能人员 & 134 & 10.42\% \\ \hline
\multicolumn{1}{|c|}{\bfseries 合计} & {\bfseries 1,286} & {\bfseries 100.00\%} \\ \hline
\end{longtable}
\end{gstab}

按学历划分，硕士及以上学历512人（占39.81\%），本科学历645人（占50.16\%），
大专及以下学历129人（占10.03\%）。
